\documentclass{Templates/customArticle}

\usepackage{float}

\begin{document}

\begin{titlepage}
    \centering
    {\LARGE\textbf{TP1 - HAI913I\\Evolution et restructuration des logiciels}\par}
    \vspace{1cm}
    {\LARGE\textbf{Analyse statique}\par}
    \vspace{1cm}
    {\Large\textbf{Nathan Boureux}\par}
    \vfill
    {\large{Master 2 Informatique, parcours Génie Logiciel}\par}
    \vspace{0.5cm}
    \includegraphics[width=0.2\textwidth]{assets/UM.png}
    \includegraphics[width=0.2\textwidth]{assets/fds.png}
\end{titlepage}

\tableofcontents
\thispagestyle{empty}
\clearpage
\pagenumbering{arabic}

\section{Partie A}

\subsection{A1}

Dans JDT, l'infrastructure visite les enfants d'un nœud si et seulement si la méthode visit a renvoyé true.

Renvoyer false élague le sous-arbre : cela signifie que le visiteur ne descendra pas dans les enfants de ce nœud.

Cependant, même si tu renvoies false, la méthode endVisit sera tout de même appelée pour ce nœud.

C'est exactement ce mécanisme qui permet, par exemple, d'ignorer volontairement le corps d'une classe anonyme lors du parcours.  

\end{document}